% ═══════════════════════════════════════════════════════════════════════
%  Macros et styles TikZ du volume — partagés par le livre (PDF) et
%  les figures autonomes converties en SVG pour le web.
%  Repris tel quel du préambule de Livre.tex.
% ═══════════════════════════════════════════════════════════════════════
% ── Bibliothèques TikZ utilisées dans les figures ────────────────────
\usetikzlibrary{calc, arrows.meta, angles, quotes, patterns, positioning,
                decorations.markings, backgrounds, fit, shapes.geometric}

% ── Raccourcis propres à ce volume ───────────────────────────────────
\newcommand{\Plan}{\Pi}                       % le plan euclidien
\newcommand{\Transf}{\mathcal{T}}             % groupe des transformations
\newcommand{\Aff}{\mathcal{A}}                % groupe des affinités
\newcommand{\Sim}{\mathcal{S}}                % groupe des similitudes
\newcommand{\Iso}{\mathcal{I}}                % groupe des isométries
\newcommand{\mat}[4]{\begin{pmatrix}#1 & #2\\ #3 & #4\end{pmatrix}}
\newcommand{\col}[2]{\begin{pmatrix}#1\\ #2\end{pmatrix}}
\newcommand{\coliii}[3]{\begin{pmatrix}#1\\ #2\\ #3\end{pmatrix}}
\newcommand{\norme}[1]{\lVert #1 \rVert}
\newcommand{\Aire}{\operatorname{Aire}}
\newcommand{\dete}{\operatorname{det}}
\renewcommand{\vec}[1]{\mathbf{#1}}

% ── Raccourcis propres au chapitre 2 (espaces métriques) ─────────────
\newcommand{\dE}{d_E}                         % distance euclidienne
\newcommand{\dT}{d_T}                         % taxi-distance (Manhattan)
\newcommand{\dS}{d_S}                         % distance sphérique (grand cercle)
\newcommand{\dHam}{d_{\mathrm{Ham}}}          % distance de Hamming (erreur)
\newcommand{\dH}{d_H}                         % distance de Hausdorff
\newcommand{\boule}[2]{B(#1, #2)}             % boule ouverte
\newcommand{\Vois}[2]{#1^{\oplus #2}}         % epsilon-voisinage de A (Minkowski)
\newcommand{\Comp}[1]{\mathcal{K}(#1)}        % ensemble des compacts non vides de E

% ── Raccourcis propres au chapitre 3 (fractales et SIF) ──────────────
\newcommand{\dims}{\dim_{\mathrm{sim}}}       % dimension de similitude
\definecolor{vertpiece}{RGB}{58, 125, 68}     % 4e couleur des « pièces »

% ── Raccourcis propres à l'annexe D (algèbre linéaire) ──────────────
\DeclareMathOperator{\Span}{span}          % espace engendré
\DeclareMathOperator{\rang}{rang}          % rang d'une matrice
\DeclareMathOperator{\Ima}{Im}             % image d'une matrice
% Vue 3D en perspective cavalière : x vers l'avant, y à droite, z en haut.
\tikzset{vue3d/.style={x={(-0.42cm,-0.32cm)}, y={(1cm,0cm)}, z={(0cm,1cm)}}}

% ── Styles TikZ communs à toutes les figures du volume ───────────────
\tikzset{
  axe/.style        = {-{Stealth[length=2mm]}, gray!70, thin},
  quadrillage/.style= {gray!22, very thin},
  fig/.style        = {fill=steelblue4!22, draw=steelblue4, thick,
                       line join=round},
  figim/.style      = {fill=firebrick3!22, draw=firebrick3, thick,
                       line join=round},
  figaux/.style     = {fill=coveraccent!22, draw=coveraccent!85!black,
                       thick, line join=round},
  pt/.style         = {circle, fill=sectioncolor, inner sep=1.1pt},
  ptim/.style       = {circle, fill=firebrick3, inner sep=1.1pt},
  fleche/.style     = {-{Stealth[length=2.4mm]}, sectioncolor, thick},
  etiq/.style       = {font=\scriptsize, inner sep=1.5pt},
  droite/.style     = {sectioncolor!70, thick, dashed},
  % Figures du chapitre 3, générées par figures/fractales/generer.py
  fractale/.style   = {fill=sectioncolor},
  nuage/.style      = {fill=sectioncolor},
  courbe/.style     = {sectioncolor, thin, line join=round},
  piece1/.style     = {fill=firebrick3},
  piece2/.style     = {fill=steelblue4},
  piece3/.style     = {fill=coveraccent},
  piece4/.style     = {fill=vertpiece},
  trait1/.style     = {firebrick3, thick, line join=round},
  trait2/.style     = {steelblue4, thick, line join=round},
  trait3/.style     = {coveraccent, thick, line join=round},
  trait4/.style     = {vertpiece, thick, line join=round},
  chaosvers/.style  = {gray!60, dashed, thin},
  chaospas/.style   = {-{Stealth[length=2.2mm]}, firebrick3, thick},
  % Chapitre 4 : états d'une chaîne de Markov, pages d'un mini-web, liens
  etat/.style       = {circle, draw=sectioncolor, thick, fill=steelblue4!10,
                       minimum size=9mm, inner sep=1pt, font=\small},
  pageweb/.style    = {etat, minimum size=7.5mm},
  lien/.style       = {-{Stealth[length=2.2mm]}, thick, sectioncolor!85!black},
  lien2/.style      = {{Stealth[length=2.2mm]}-{Stealth[length=2.2mm]}, thick,
                       sectioncolor!85!black},
  proba/.style      = {font=\scriptsize, fill=white, inner sep=1pt},
}

% Repère quadrillé : \repere{xmin}{xmax}{ymin}{ymax}
\newcommand{\repere}[4]{%
  \draw[quadrillage] (#1,#3) grid (#2,#4);
  \draw[axe] (#1,0) -- (#2+0.35,0) node[etiq, below] {$x$};
  \draw[axe] (0,#3) -- (0,#4+0.35) node[etiq, left] {$y$};}

% Le motif « F », tracé dans le carré unité : sert à voir d'un coup d'œil
% si une transformation retourne le plan ou non.
\newcommand{\motifF}{(0,0) -- (0.3,0) -- (0.3,0.55) -- (0.6,0.55) --
  (0.6,0.75) -- (0.3,0.75) -- (0.3,0.85) -- (0.75,0.85) -- (0.75,1) --
  (0,1) -- cycle}
